% GENERATED FILE. Edit canonical lessons, then run npm run generate:books.
% canonical-source-hash: fnv1a64:4c08b9db01eee708
% canonical-lessons: PT-C312-noivado, PT-C312-alianca, PT-C312-casamento, PT-C312-pedir-em-casamento, PT-C312-matrimonio

\chapter{Noivos}
\label{ch:pt-noivos}
\hlchaptermodality{312}

\begin{chapteropening}
\textbf{By the end of this chapter:} \emph{I can talk about getting engaged, rings, weddings, proposals and marriage.}

Complete the last lesson of chapter 312: I can talk about getting engaged, rings, weddings, proposals and marriage.
\end{chapteropening}

\section[o noivado]{o noivado — the engagement}

\label{lesson:PT-C312-noivado}



\begin{quote}
Before the new one: say the Portuguese for old age, then the Portuguese for the generation.
\end{quote}

\subsection*{You'll want to know: o noivado}
\textbf{o noivado} — \textquotedblleft{}the engagement\textquotedblright{}.

\begin{quote}
O noivado durou dois anos.
\end{quote}

\begin{quote}
\emph{The engagement lasted two years.}
\end{quote}

Built on \emph{noivo}, ‘fiancé, groom’.

\begin{cousinweb}
\emph{noivado} is built on \emph{noivo} (EtymDB), which goes back to Vulgar Latin \emph{novius} \textquotedblleft{}newly wed\textquotedblright{}, from Latin \emph{novus} \textquotedblleft{}new\textquotedblright{}.
\end{cousinweb}

\begin{grammarlens}[title={What you've built}]
One more word for telling, explaining and describing at B1.
\end{grammarlens}

\subsection*{Your turn}
\begin{itemize}
\raggedright
  \item \emph{Say it:} \emph{o noivado}
  \item \emph{Say it:} \emph{o noivado}, once more
  \item \emph{Say it:} a sentence of your own with \emph{o noivado}
  \item \emph{Recall:} say the Portuguese for old age, then the Portuguese for the generation, then say \emph{o noivado} again
\end{itemize}

\begin{tcolorbox}[breakable,colback=teal!4,colframe=teal!35!black,title={Before you move on}]
What does o noivado mean? (\textquotedblleft{}The engagement\textquotedblright{}.) Say it once more.
\end{tcolorbox}
\section[a aliança]{a aliança — the wedding ring; the alliance}

\label{lesson:PT-C312-alianca}



\begin{quote}
Before the new one: say the Portuguese for the generation, then the Portuguese for the engagement.
\end{quote}

\subsection*{You'll want to know: a aliança}
\textbf{a aliança} — \textquotedblleft{}the wedding ring; the alliance\textquotedblright{}.

\begin{quote}
Ele usa a aliança na mão esquerda.
\end{quote}

\begin{quote}
\emph{He wears his wedding ring on his left hand.}
\end{quote}

Also a political alliance.

\begin{cousinweb}
\emph{aliança} came from French \emph{alliance} \textquotedblleft{}a union, a marriage\textquotedblright{} (EtymDB), built on \emph{allier} \textquotedblleft{}to bind\textquotedblright{}, from Latin \emph{alligāre} \textquotedblleft{}to tie to\textquotedblright{} (Lewis \& Short).
\end{cousinweb}

\begin{grammarlens}[title={What you've built}]
One more word for telling, explaining and describing at B1.
\end{grammarlens}

\subsection*{Your turn}
\begin{itemize}
\raggedright
  \item \emph{Say it:} \emph{a aliança}
  \item \emph{Say it:} \emph{a aliança}, once more
  \item \emph{Say it:} a sentence of your own with \emph{a aliança}
  \item \emph{Recall:} say the Portuguese for the generation, then the Portuguese for the engagement, then say \emph{a aliança} again
\end{itemize}

\begin{tcolorbox}[breakable,colback=teal!4,colframe=teal!35!black,title={Before you move on}]
What does a aliança mean? (\textquotedblleft{}The wedding ring; the alliance\textquotedblright{}.) Say it once more.
\end{tcolorbox}
\section[o casamento]{o casamento — the wedding; marriage}

\label{lesson:PT-C312-casamento}



\begin{quote}
Before the new one: say the Portuguese for the engagement, then the Portuguese for the wedding ring.
\end{quote}

\subsection*{You'll want to know: o casamento}
\textbf{o casamento} — \textquotedblleft{}the wedding; marriage\textquotedblright{}.

\begin{quote}
O casamento foi numa quinta perto do rio.
\end{quote}

\begin{quote}
\emph{The wedding was at a country estate near the river.}
\end{quote}

Built on \emph{casar}.

\begin{cousinweb}
\emph{casamento} comes down from Old Portuguese (EtymDB), built on \emph{casar} \textquotedblleft{}to marry\textquotedblright{}, from \emph{casa} \textquotedblleft{}a house\textquotedblright{}, which is inherited from Latin \emph{casa} \textquotedblleft{}a hut, a cottage\textquotedblright{}: to set up house.
\end{cousinweb}

\begin{grammarlens}[title={What you've built}]
One more word for telling, explaining and describing at B1.
\end{grammarlens}

\subsection*{Your turn}
\begin{itemize}
\raggedright
  \item \emph{Say it:} \emph{o casamento}
  \item \emph{Say it:} \emph{o casamento}, once more
  \item \emph{Say it:} a sentence of your own with \emph{o casamento}
  \item \emph{Recall:} say the Portuguese for the engagement, then the Portuguese for the wedding ring, then say \emph{o casamento} again
\end{itemize}

\begin{tcolorbox}[breakable,colback=teal!4,colframe=teal!35!black,title={Before you move on}]
What does o casamento mean? (\textquotedblleft{}The wedding; marriage\textquotedblright{}.) Say it once more.
\end{tcolorbox}
\section[pedir em casamento]{pedir em casamento — to propose (marriage)}

\label{lesson:PT-C312-pedir-em-casamento}



\begin{quote}
Before the new one: say the Portuguese for the wedding ring, then the Portuguese for the wedding.
\end{quote}

\subsection*{You'll want to know: pedir em casamento}
\textbf{pedir em casamento} — \textquotedblleft{}to propose (marriage)\textquotedblright{}.

\begin{quote}
Ele pediu-a em casamento na praia.
\end{quote}

\begin{quote}
\emph{He proposed to her on the beach.}
\end{quote}

\emph{Pedir alguém em casamento}; the pronoun goes on \emph{pedir}: \emph{pediu-a}.

\begin{cousinweb}
\emph{pedir} is inherited from Latin \emph{petere} \textquotedblleft{}to seek, to ask for\textquotedblright{} (Lewis \& Short); \emph{casamento} is built on \emph{casar}, from \emph{casa} \textquotedblleft{}a house\textquotedblright{}.
\end{cousinweb}

\begin{grammarlens}[title={What you've built}]
One more word for telling, explaining and describing at B1.
\end{grammarlens}

\subsection*{Your turn}
\begin{itemize}
\raggedright
  \item \emph{Say it:} \emph{pedir em casamento}
  \item \emph{Say it:} \emph{pedir em casamento}, once more
  \item \emph{Say it:} a sentence of your own with \emph{pedir em casamento}
  \item \emph{Recall:} say the Portuguese for the wedding ring, then the Portuguese for the wedding, then say \emph{pedir em casamento} again
\end{itemize}

\begin{tcolorbox}[breakable,colback=teal!4,colframe=teal!35!black,title={Before you move on}]
What does pedir em casamento mean? (\textquotedblleft{}To propose (marriage)\textquotedblright{}.) Say it once more.
\end{tcolorbox}
\section[o matrimónio]{o matrimónio — matrimony, marriage (formal)}

\label{lesson:PT-C312-matrimonio}



\begin{quote}
Before the new one: say the Portuguese for the wedding, then the Portuguese for to propose.
\end{quote}

\subsection*{You'll want to know: o matrimónio}
\textbf{o matrimónio} — \textquotedblleft{}matrimony, marriage (formal)\textquotedblright{}.

\begin{quote}
Celebraram o matrimónio numa pequena igreja.
\end{quote}

\begin{quote}
\emph{They celebrated their marriage in a small church.}
\end{quote}

In Brazil: \emph{matrimônio}.

\begin{cousinweb}
\emph{matrimónio} is a learned borrowing from Latin \emph{mātrimōnium} \textquotedblleft{}marriage\textquotedblright{} (EtymDB), built on \emph{māter} \textquotedblleft{}mother\textquotedblright{}.
\end{cousinweb}

\begin{grammarlens}[title={What you've built}]
One more word for telling, explaining and describing at B1.
\end{grammarlens}

\subsection*{Your turn}
\begin{itemize}
\raggedright
  \item \emph{Say it:} \emph{o matrimónio}
  \item \emph{Say it:} \emph{o matrimónio}, once more
  \item \emph{Say it:} a sentence of your own with \emph{o matrimónio}
  \item \emph{Recall:} say the Portuguese for the wedding, then the Portuguese for to propose, then say \emph{o matrimónio} again
\end{itemize}

\begin{tcolorbox}[breakable,colback=teal!4,colframe=teal!35!black,title={Before you move on}]
What does o matrimónio mean? (\textquotedblleft{}Matrimony, marriage (formal)\textquotedblright{}.) Say it once more.
\end{tcolorbox}
